\documentclass[11pt]{article}

\usepackage[hmargin=0.82in,vmargin=0.72in]{geometry}
\usepackage[T1]{fontenc}
\usepackage{lmodern}
\usepackage{microtype}
\usepackage{amsmath,amssymb,amsthm}
\usepackage{booktabs}
\usepackage{enumitem}
\usepackage{titlesec}
\usepackage{caption}
\usepackage{pgfplots}
\usepackage[hidelinks]{hyperref}

\pgfplotsset{compat=1.17}
\setlength{\parindent}{0pt}
\setlength{\parskip}{4pt}
\setlist[itemize]{leftmargin=1.4em,itemsep=2pt,topsep=3pt}
\titlespacing*{\section}{0pt}{9pt}{3pt}
\titlespacing*{\subsection}{0pt}{7pt}{2pt}
\titleformat{\section}{\large\bfseries}{}{0pt}{}
\titleformat{\subsection}{\normalsize\bfseries}{}{0pt}{}
\captionsetup{font=small,labelfont=bf,skip=4pt}
\newtheorem{proposition}{Proposition}

\newcommand{\dY}{d_Y}
\newcommand{\dO}{d_O}

\begin{document}

\begin{center}
{\LARGE\bfseries Demographic Momentum and Intergenerational Housing}\par
\vspace{3pt}
{\large A two-generation extension of reproductive equilibrium}\par
\vspace{3pt}
{\small August 12, 2026}
\end{center}

\vspace{-2pt}
Provided a new positive steady state exists and is stable, the steady-state
theory says that a permanent decline in the desire for children eventually
produces a smaller adult-household population and a lower house price, with
fertility returning to effective replacement. The apparent difficulty is
the transition: U.S. population has continued to grow, and house prices can
rise while fertility is below replacement. This note shows exactly when those
facts are consistent with the theory.

The missing state variable is cohort composition. A large young cohort can
produce fewer children than replacement and still be larger than the old cohort
it replaces. As that cohort ages, its housing demand can also outweigh the loss
of housing demand from the smaller cohort behind it. Population and prices may
therefore rise temporarily even though the economy is moving toward a smaller
steady state. This mechanism connects the reproductive closure to the paper's
intergenerational housing channel without replacing its lifecycle model.

\section*{1. A two-generation economy}

At date $t$, $Y_t$ is the mass of young-family households and $O_t$ is the mass
of old-family households. A young family chooses consumption $c$, adult housing $h$, and
children $n$:
\begin{equation}
 \max_{c,h,n}\;c+\alpha\log h+\theta\log n
 \quad\text{subject to}\quad
 c+ph+\bigl[q+a(p+z)\bigr]n=y .
 \label{eq:young_problem}
\end{equation}
The market cost of housing is $p$. Each child uses $a$ units of space and costs
$q$ units of other goods. The term $z\geq0$ is a goods-equivalent shadow cost
of obtaining child space. In the full paper it summarizes rental size limits,
down-payment constraints, and other access wedges faced by young families.
The interior demands are
\begin{equation}
 h_Y(p)=\frac{\alpha}{p},
 \qquad
 n(p;\theta,z)=\frac{\theta}{q+a(p+z)},
 \qquad
 \dY(p;\theta,z)=\frac{\alpha}{p}+a n(p;\theta,z).
 \label{eq:young_demands}
\end{equation}
Here $\dY$ is a young family's total housing demand, including space used by
children. We assume income is high enough for the interior solution.

An old family demands
\begin{equation}
 \dO(p;\tau)=\frac{\beta(\tau)}{p},
 \qquad \beta'(\tau)>0.
 \label{eq:old_demand}
\end{equation}
This is a reduced-form summary of the old household's tenure and retention
problem in the main paper. A larger $\tau$ makes old households retain more
housing because moving is costly or housing is valued as an estate asset.
Housing supply is $H^s(p)=A p^\varepsilon$, with $A>0$ and $\varepsilon\geq0$.
The market-clearing price satisfies
\begin{equation}
 A p_t^\varepsilon
 =Y_t\dY(p_t;\theta,z)+O_t\dO(p_t;\tau).
 \label{eq:clearing}
\end{equation}

Let $\bar n$ denote effective replacement fertility. The timing is deliberately
simple: children chosen at $t$ form the next young generation, while the current
young become old. Thus
\begin{equation}
 Y_{t+1}=R(p_t;\theta,z)Y_t,
 \qquad O_{t+1}=Y_t,
 \qquad
 R(p;\theta,z)\equiv\frac{n(p;\theta,z)}{\bar n}.
 \label{eq:cohort_law}
\end{equation}
The two equations in~\eqref{eq:cohort_law} are the only addition to the earlier
steady-state sketch.

\section*{2. The new steady state}

A positive stationary population requires $Y^*=O^*$ and $R(p^*)=1$.
Reproduction therefore selects the house price; housing clearing selects the
number of families at that price.

\begin{proposition}[Reproductive steady state]
\label{prop:ss}
A unique positive steady state exists if
$\theta>\bar n(q+az)$. It is
\begin{equation}
 p^*=\frac{\theta/\bar n-q}{a}-z,
 \qquad
 Y^*=O^*
 =\frac{A(p^*)^{\varepsilon+1}}
 {\alpha+\beta(\tau)+a\bar n p^*},
 \qquad N^*=2Y^* .
 \label{eq:ss}
\end{equation}
A lower $\theta$ or a higher young access wedge $z$ lowers both $p^*$ and
$N^*$. Greater housing retention by the old raises $\beta(\tau)$ and lowers
$N^*$, but does not change the reproductive price $p^*$.
\end{proposition}

The last result is the intergenerational counterpart of demographic
capitalization. At a reproductive steady state, old households cannot
permanently raise the price: a higher price would push fertility below
replacement. Instead, stronger old-age housing retention reduces the number of
families that the housing stock can support at the demographically pinned
price. Conversely, a policy that eases young access ($z\downarrow$) raises the
long-run price and population, while a policy that releases old housing
($\tau\downarrow$) raises population without changing the long-run price.
Both policies leave effective reproduction equal to one in the new steady
state; they change the scale of the economy, not its stationary growth rate.

\section*{3. When can household population and prices still rise?}

Fix the post-shock value of $\theta$ and define
$r_t=R(p_t;\theta,z)$ and $\mu_t=O_t/Y_t$. Suppose fertility is below
replacement at date $t$, so $r_t<1$. The mass of young families is shrinking,
but total adult-household population need not shrink: the relevant comparison is between
the new young cohort and the old cohort that exits.

\begin{proposition}[Demographic momentum]
\label{prop:momentum}
With $r_t<1$, total adult-household population rises from $t$ to $t+1$ if and only if
\begin{equation}
 \mu_t<r_t<1.
 \label{eq:population_growth}
\end{equation}
The house price rises if and only if
\begin{equation}
 (1-\mu_t)\dO(p_t;\tau)
 >(1-r_t)\dY(p_t;\theta,z).
 \label{eq:price_growth}
\end{equation}
Hence below-replacement fertility, adult-household population growth, and house-price growth
coexist if and only if
\begin{equation}
 \mu_t<\min\left\{
 r_t,\;
 1-(1-r_t)\frac{\dY(p_t;\theta,z)}{\dO(p_t;\tau)}
 \right\}.
 \label{eq:joint_condition}
\end{equation}
\end{proposition}

Equation~\eqref{eq:price_growth} separates two forces. The smaller next young
cohort removes $(1-r_t)Y_t\dY$ units of demand. The unusually large current
young cohort becomes old and adds $(Y_t-O_t)\dO$ units relative to the cohort
that exits. Prices rise when the second force is larger.

This proposition also clarifies where the transition can start. If the economy
begins from an \emph{exact} stationary age distribution, then $\mu_t=1$.
Below-replacement fertility immediately makes both population and the price
fall. The proposed interpretation therefore cannot be ``the economy was
exactly in steady state and then fertility fell.'' It can be: aggregate
population and prices were near their old steady-state values, but a large
young cohort was moving through the age distribution. That cohort bulge may
reflect earlier fertility, immigration, or both.

\section*{4. A numerical illustration}

Figure~\ref{fig:transition} uses a fixed housing stock and normalizes effective
replacement to one. The initial economy has the old steady-state household mass
and price, but its composition is tilted toward the young. The illustrative
parameters make young and old housing demand equal at the old price, so this
nonstationary composition preserves both aggregates exactly. At date zero,
$\theta$ falls permanently. The impact price falls to $0.941$ and effective
reproduction falls to $0.888$. In the next period, however, the initial young
bulge becomes old: household mass rises by $8.7$ percent and the price rises to
$1.051$, even though reproduction remains below replacement. The economy then
converges with a cohort echo to a new steady state with price $0.700$ and
household mass $22.8$ percent below the old steady state.

\begin{figure}[ht!]
\centering
\begin{tikzpicture}
\begin{axis}[
 name=price,
 width=0.43\textwidth,
 height=0.31\textwidth,
 xmin=0,xmax=10,
 ymin=0.64,ymax=1.23,
 xtick={0,2,4,6,8,10},
 xlabel={Generation after the shock},
 ylabel={Price / old steady state},
 title={(a) House price},
 axis lines=left,
 tick label style={font=\small},
 label style={font=\small},
 title style={font=\small\bfseries},
 legend style={font=\footnotesize,draw=none,fill=none,at={(0.98,0.98)},anchor=north east},
]
\addplot[very thick,red!75!black,mark=*] table[
 x=period,y=price_shock,col sep=comma
] {../output/model/demographic_momentum_intergenerational_housing_note/transition_paths.csv};
\addlegendentry{Fertility shock}
\addplot[thick,blue!70!black,dashed] table[
 x=period,y=price_no_shock,col sep=comma
] {../output/model/demographic_momentum_intergenerational_housing_note/transition_paths.csv};
\addlegendentry{No-shock path}
\addplot[thin,black,dotted] coordinates {(0,1) (10,1)};
\addplot[thin,red!75!black,dotted] coordinates {(0,0.7) (10,0.7)};
\end{axis}

\begin{axis}[
 at={(price.east)},anchor=west,xshift=0.105\textwidth,
 width=0.43\textwidth,
 height=0.31\textwidth,
 xmin=0,xmax=10,
 ymin=0.70,ymax=1.19,
 xtick={0,2,4,6,8,10},
 xlabel={Generation after the shock},
 ylabel={Household mass / old steady state},
 title={(b) Adult-household mass},
 axis lines=left,
 tick label style={font=\small},
 label style={font=\small},
 title style={font=\small\bfseries},
]
\addplot[very thick,red!75!black,mark=*] table[
 x=period,y=population_shock,col sep=comma
] {../output/model/demographic_momentum_intergenerational_housing_note/transition_paths.csv};
\addplot[thick,blue!70!black,dashed] table[
 x=period,y=population_no_shock,col sep=comma
] {../output/model/demographic_momentum_intergenerational_housing_note/transition_paths.csv};
\addplot[thin,black,dotted] coordinates {(0,1) (10,1)};
\addplot[thin,red!75!black,dotted] coordinates {(0,0.7724138) (10,0.7724138)};
\end{axis}
\end{tikzpicture}
\caption{A fertility decline during a cohort transition. Solid lines impose a
permanent decline in the value of children; dashed lines retain its old value.
Black dotted lines are the old steady state and red dotted lines are the new
steady state. The initial young-cohort bulge creates temporary population and
price growth, but the shock path lies below the no-shock counterfactual. The
parameters are illustrative, not calibrated.}
\label{fig:transition}
\end{figure}

The counterfactual in Figure~\ref{fig:transition} is important. The preference
shock does not \emph{cause} the price increase; it shifts young families'
willingness to pay inward and lowers the clearing price relative to the
no-shock path in this illustration. Cohort composition can
nevertheless make the observed price rise over time. A credible empirical
application must therefore discipline the initial age distribution and compare
paths, rather than infer the sign of the shock from the sign of price growth.

\section*{5. How this enters the full paper}

The two-generation model should remain a theory layer, not replace the
quantitative household problem. The full model already contains young access
wedges, tenure choice, old-age housing retention, fertility, survival, and the
housing market. What it lacks for this exercise is a calendar-time law for
\emph{unnormalized} cohort masses. In general notation, let $m_t$ be the measure
of households over age, wealth, tenure, location, and children. The extension is
\begin{equation}
 m_{t+1}=\mathcal T(p_t;\theta_t,\pi)m_t+\iota_{t+1},
 \qquad
 H^s(p_t)=\int d(s,p_t;\theta_t,\pi)\,dm_t(s),
 \label{eq:full_transition}
\end{equation}
where $\mathcal T$ contains survival, household choices, and the conversion of
locally born children into young entrants; $\iota_{t+1}$ is the measure of
outside entrants, including immigration. The policy vector $\pi$ contains the
paper's access and retention institutions. Unlike a stationary computation,
equation~\eqref{eq:full_transition} does not renormalize the population after
each iteration.

This structure preserves the paper's intergenerational focus in two ways.
First, tenure and transaction costs determine how much housing the older cohort
retains, which changes prices faced by the young and therefore fertility.
Second, policies aimed at young access change both today's tenure allocation
and the size of future cohorts. The same model can then report an impact effect,
a cohort transition, and a new reproductive steady state.

The disciplined next step is modest: initialize $m_0$ from the joint
distribution of age, wealth, tenure, location, and children (using the model
where the full joint distribution is unavailable); specify transparent paths
for the child-preference shifter and immigration; and ask whether the model can
reproduce the joint path of cohort size, fertility, and housing prices. This
note establishes the possibility and the necessary inequalities. It does not
claim that U.S. cohort composition is quantitatively sufficient. Nor does the
generic construction guarantee that the maintained quantitative
parameterization has a positive reproductive root: the latest closed-economy
audit does not find one over the tested price range. Establishing such a root
and its stability is a prerequisite for using this as the quantitative closure.

\appendix
\section*{Appendix: proofs and local dynamics}

\textbf{Proof of Proposition~\ref{prop:ss}.}
At a positive steady state, equation~\eqref{eq:cohort_law} requires
$R(p^*)=1$. Solving $n(p^*)=\bar n$ gives the price in~\eqref{eq:ss}; it is
positive exactly when $\theta>\bar n(q+az)$. Since $Y^*=O^*$, multiply housing
clearing by $p^*$ and substitute $n(p^*)=\bar n$ to obtain the expression for
$Y^*$. Moreover,
\begin{equation}
 \frac{dY^*}{dp^*}
 =\frac{A(p^*)^\varepsilon
 \left[(\varepsilon+1)(\alpha+\beta)
 +\varepsilon a\bar n p^*\right]}
 {\left[\alpha+\beta+a\bar n p^*\right]^2}>0.
\end{equation}
The comparative statics follow from $p^*_{\theta}=1/(a\bar n)>0$,
$p^*_{z}=-1$, and $Y^*_{\beta}<0$ holding $p^*$ fixed. \hfill$\square$

\textbf{Proof of Proposition~\ref{prop:momentum}.}
Total population is $N_t=(1+\mu_t)Y_t$, while
$N_{t+1}=(1+r_t)Y_t$. This proves~\eqref{eq:population_growth}. Evaluate next
period's housing demand at the current price and subtract current demand:
\begin{align}
 &Y_{t+1}\dY(p_t)+O_{t+1}\dO(p_t)
 -Y_t\dY(p_t)-O_t\dO(p_t) \notag\\
 &\hspace{2cm}=Y_t\left[(r_t-1)\dY(p_t)
 +(1-\mu_t)\dO(p_t)\right].
\end{align}
Supply is weakly increasing and demand is strictly decreasing in $p$, so the
sign of this expression is the sign of $p_{t+1}-p_t$. This proves
\eqref{eq:price_growth}; combining the two inequalities gives
\eqref{eq:joint_condition}. \hfill$\square$

\textbf{Local cohort echoes.}
Let $f_t=\log(Y_t/Y^*)$. Write $\eta_Y>0$ and $\eta_O>0$ for the elasticities of
the market-clearing price with respect to young and old cohort masses, and let
$\epsilon_R<0$ be the price elasticity of effective reproduction, all evaluated
at the steady state. Since $O_t=Y_{t-1}$, a first-order approximation is
\begin{equation}
 f_{t+1}=(1+\epsilon_R\eta_Y)f_t
 +\epsilon_R\eta_O f_{t-1}.
 \label{eq:local_law}
\end{equation}
Set $a_R=-\epsilon_R\eta_Y>0$ and $b_R=-\epsilon_R\eta_O>0$. The characteristic
polynomial is $\lambda^2-(1-a_R)\lambda+b_R$. The steady state is locally stable
if and only if $b_R<1$ and $a_R-b_R<2$. Convergence is oscillatory when
$(1-a_R)^2<4b_R$. Old-age housing demand creates the lagged term $b_R$ and is
therefore the source of the cohort echo visible in Figure~\ref{fig:transition}.

\end{document}
